\documentclass[10pt,a4paper]{amsart}
\usepackage[margin=19mm]{geometry}
\usepackage{amsmath,amssymb,mathtools}
\usepackage[scr=rsfs,cal=euler]{mathalfa}
\usepackage{xfrac}
\usepackage[unicode,hidelinks,bookmarks=false]{hyperref}
\newcommand{\ie}{i.e.\ }
\newcommand{\cf}{cf.\ }
\newcommand{\resp}{respectively\ }
\newcommand{\Subsection}{\subsection}
\providecommand{\nobreakdash}{\nobreak\hbox{-}}
\setlength{\emergencystretch}{2em}
\multlinegap0pt
\usepackage{siunitx}
\usepackage{siunitx}
\usepackage{siunitx}
\usepackage{siunitx}
\usepackage{amssymb}
\usepackage{siunitx}
\usepackage{float}
\usepackage{algorithm}\usepackage{algpseudocode}
\usepackage{xcolor}
\newtheorem{theorem}{Theorem}
\newtheorem{proposition}{Proposition}
\usepackage{cleveref}
\crefname{theorem}{Theorem}{Theorems}
\crefname{proposition}{Proposition}{Propositions}
\newcommand\calG{\mathcal{G}}
\newcommand\calR{\mathcal{R}}
\title{Data-informed posterior approximation for Bayesian linear inverse problems}
\author{Haibo Li}
\date{Source: arXiv:2605.20770v1 (CC BY 4.0)}
\begin{document}
\maketitle

\noindent\emph{Source and licence.} Data-informed posterior approximation for Bayesian linear inverse problems. Version arXiv:2605.20770v1 (CC BY 4.0), distributed by arXiv under the Creative Commons Attribution 4.0 International licence, http://creativecommons.org/licenses/by/4.0/, as recorded in the arXiv licence metadata for this exact version and shown on the abstract page https://arxiv.org/abs/2605.20770v1. Full licence notice: \texttt{licenses/arxiv-2605.20770-CC-BY-4.0.txt}.\par\smallskip

\noindent\textit{Editorial scope.} Counted unit: the article's theorem thm:qgkb-lanczos (source0.tex lines 552--558). The article's complete original proof of it is delivered with it (source0.tex lines 559--597, one \texttt{proof} environment of 39 lines). No mathematics is written, repaired or re-derived: the statement and the proof are the article's own lines, in the article's own notation, and no retained line is substituted or re-flowed.  Retained uncounted as the source-local context the proof needs: lines 263--291; lines 472--490.  Nothing was omitted from any retained span, so the package carries no omission record.  No retained line contains a citation, so the wrapper carries no bibliography and no bibliography entry.  No retained line contains a figure environment, an image-inclusion command or drawing code, so no image is delivered and the package declares no asset dependency.  Editorial formatting only, in addition: an AMS \texttt{amsart} preamble that restates the article's own declarations and macros used by the retained lines, namely the environments \texttt{theorem}, \texttt{proposition} on separate counters, together with the standard packages those lines need.  The retained environments print in this document as Theorem 1, Proposition 1 in order of appearance, whereas the article prints the same items with the numbering of its own full document.  The complete native source of the article as distributed in the arXiv e-print for this exact version is bundled unchanged as \texttt{sources/201016/source0.tex}.
\begin{theorem}[Correspondence between data and parameter spaces]\label{thm:isom}
% \todo[color=green!20]{Haibo: verify theorem}
  Define the linear operator 
  \begin{align}\label{isom}
    \begin{split}
      \iota: (\mathbb R^m/\mathcal N(M),\langle\cdot,\cdot\rangle_M)
          &\longrightarrow (\mathbb R^n,\langle\cdot,\cdot\rangle_{\Sigma^{-1}}) \\
      [u] &\longmapsto \Sigma G^{\top}u.
    \end{split}
  \end{align} 
  Then the following statements hold.
\begin{enumerate}
\item[(a)] (Isometric embedding)
The map $\iota$ is an isometric embedding, i.e.,
\[ \langle \iota([u]),\iota([v])\rangle_{\Sigma^{-1}}
=\langle [u],[v]\rangle_M,
\]
and $\iota$ is injective.
\item[(b)] (Parameter space decomposition)
It holds that 
\begin{equation*}
  (\mathbb R^n,\langle\cdot,\cdot\rangle_{\Sigma^{-1}})
  = \mathrm{Im}(\iota) \oplus \mathcal{N}(G) ,
\end{equation*}
where the two subspaces are $\Sigma^{-1}$-orthogonal, and $x_{\lambda}\in \mathrm{Im}(\iota)=\mathcal{R}(\Sigma G^\top)$.
\item[(c)] (Spectral correspondence)
Let $\{(\hat\mu_i,\hat w_i)\}_{i=1}^{\hat{l}}$ be all nonzero $M$-orthonormal generalized eigenpairs of $(M,\,\Gamma)$. Then $l=\hat{l}=\mathrm{rank}(G)$, and $\{\hat\mu_i,\iota([\hat w_i])\}_{i=1}^{l}$ are all nonzero $\Sigma^{-1}$-orthonormal eigenpairs of $(H,\,\Sigma^{-1})$.
\end{enumerate}
\end{theorem}

\begin{proposition}\label{thm:qsggkb-properties}
Assume that Q-GKB runs for $k$ steps with $k\leq k_b$. Then the following properties hold.
\begin{enumerate}
\item[(a)] The matrix-form relations hold:
\begin{equation}\label{eq:bidiag-relations}
\left\{
\begin{aligned}
& \beta_1 U_{k+1} e_1 = y, \\
& M V_{k} = U_{k+1} B_k, \\
& \Gamma^{-1} U_{k+1} = V_{k} B_k^\top + \alpha_{k+1} v_{k+1} e_{k+1}^\top .
\end{aligned}
\right.
\end{equation}
\item[(b)] The vectors $\{u_i\}_{i=1}^k$ form an $\Gamma^{-1}$-orthonormal basis of the Krylov subspace
\[\mathcal K_k(M\Gamma^{-1},y) = \operatorname{span}\{(M\Gamma^{-1})^i y\}_{i=0}^{k-1},\]
and vectors $\{v_i\}_{i=1}^k$ form an $M$-orthonormal basis of the Krylov subspace
\[\mathcal K_k(\Gamma^{-1}M,\Gamma^{-1}y) = \operatorname{span}\{(\Gamma^{-1}M)^i\Gamma^{-1}y\}_{i=0}^{k-1}.\]
\end{enumerate}
\end{proposition}

\begin{theorem}\label{thm:qgkb-lanczos}
  Assume that Q-GKB runs for $k$ steps with $k\leq k_b$, and the nonzero 2-orthonormal eigenpairs of $B_{k}^{\top}B_{k}$ are $\left\{(\theta_{i}^{(k)},\,s_{i}^{(k)})\right\}_{i=1}^{k}$. Then the following statements hold.
\begin{enumerate}
  \item[(a)] The pair $\left(\theta_{i}^{(k)}, p_{i}^{(k)}\right):=\left(\theta_{i}^{(k)},\iota\left([V_{k}s_{i}^{(k)}]\right)\right)$ gradually approximates a nonzero generalized eigenpair of $(H,\,\Sigma^{-1})$ as $k$ increases, and the approximation becomes exact at $k=k_b$.
  \item[(b)] At the breakdown step, $\{p_{i}^{(k_b)}\}_{i=1}^{k_b}$ are $\Sigma^{-1}$-orthonormal generalized eigenvectors of $(H,\,\Sigma^{-1})$. Moreover, each $p_i^{(k_b)}$ is the image under $\iota$ of a generalized eigenvector in one of those $\mathcal G_j$ for which $P_{\mathcal G_j}\Gamma^{-1}y\neq 0$. 
\end{enumerate}
\end{theorem}

\begin{proof}
  (a) Using the matrix form relations of Q-GKB, we have
  \begin{align*}
    \Gamma^{-1}MV_{k} = \Gamma^{-1}U_{k+1}B_{k}
    = (V_{k}B_{k}^{\top}+\alpha_{k+1}v_{k+1}e_{k+1}^{\top})B_{k}.
  \end{align*}
  A simple calculation leads to 
  \begin{align*}
    MV_{k}s_{i}^{(k)} - \theta_{i}^{(k)}\Gamma V_{k}s_{i}^{(k)}
    &= \Gamma(V_{k}B_{k}^{\top}+\alpha_{k+1}v_{k+1}e_{k+1}^{\top})B_{k}s_{i}^{(k)} - \theta_{i}^{(k)}\Gamma V_{k}s_{i}^{(k)} \\
    &= \Gamma V_{k}(B_{k}^{\top}B_{k}s_{i}^{(k)}-\theta_{i}^{(k)}s_{i}^{(k)}) + \alpha_{k+1}\beta_{k+1}\Gamma v_{k+1}e_{k}^{\top}s_{i}^{(k)} \\
    &= \alpha_{k+1}\beta_{k+1}\Gamma v_{k+1}e_{k}^{\top}s_{i}^{(k)}.
  \end{align*}
  Therefore, the pair $\left(\theta_{i}^{(k)}, V_{k}s_{i}^{(k)}\right)$ gradually approximates a nonzero generalized eigenpair of $(M,\Gamma)$ as $k$ increases, and at the breakdown step, $\left(\theta_{i}^{(k_b)}, V_{k}s_{i}^{(k_b)}\right)$ is an exact eigenpair since $\alpha_{k_b+1}\beta_{k_b+1}=0$. By \Cref{thm:isom}, $\left(\theta_{i}^{(k)},\iota\left([V_{k}s_{i}^{(k)}]\right)\right)$ gradually approximates a nonzero eigenpair of $(H, \Sigma^{-1})$, which becomes an exact one at the breakdown step.

  (b) Without loss of generality, we suppose that only the first $q$ elements in $\{P_{\mathcal{G}_1}\Gamma^{-1}y,\dots,P_{\mathcal{G}_t}\Gamma^{-1}y\}$ are nonzero. Noticing that $\{V_{k_b}s_{i}^{(k_b)}\}_{i=1}^{k_b}$ are mutually $M$-orthonormal, by \Cref{thm:isom} (a) and (c), we only need to prove that the $k_b$ vectors $\hat{p}_{i}^{(k_b)}:=V_{k_b}s_{i}^{(k)}$ belong separately to the $M$-orthogonal subspaces $\mathcal{G}_{1},\dots,\mathcal{G}_q$. Since $\theta_{i}^{(k_b)}>0$ have different values and $\calG_i$ are mutually $M$-orthogonal, these $\hat{p}_{i}^{(k_b)}$ must belong to different subspaces among $\{\mathcal{G}_i\}_{i=1}^{t}$. Therefore, we only need to prove $P_{\calG}\hat{p}_{i}^{(k_b)}=\hat{p}_{i}^{(k_b)}$ for each $1\leq i\leq q$, where $\calG=\calG_1\oplus\cdots\oplus\calG_q$ and $P_{\calG}$ is the $M$-orthogonal projector onto $\calG$. Suppose $W_i$ is an $M$-orthonormal matrix whose columns span $\mathcal{G}_i$. Then $MW_i=\mu_{i}\Gamma W_i$ and $P_{\mathcal{G}_i}=W_{i}W_{i}^{\top}M$, which leads to 
  \[\Gamma^{-1}MP_{\mathcal{G}_i}=\Gamma^{-1}MW_{i}W_{i}^{\top}M=\mu_i W_{i}W_{i}^{\top}M=\mu_iP_{\mathcal{G}_i}.\] 
  Therefore, it holds that 
  \[\Gamma^{-1}Mx 
    = \Gamma^{-1}M\sum_{i=1}^{t}P_{\mathcal{G}_i}x
    = \sum_{i=1}^{t}\mu_{i}P_{\mathcal{G}_i}x, \quad  \forall \, x\in\mathbb{R}^{m} .\]
  For any $j\geq 0$, it follows that 
  \begin{equation*}
    \widetilde{w}_i:=(\Gamma^{-1}M)^{i}\Gamma^{-1}y = \left(\sum_{i=1}^{t}\mu_{i}P_{\mathcal{G}_i} \right)^i \Gamma^{-1}y
    = \sum_{i=1}^{t}\mu_{i}^{j}P_{\mathcal{G}_i}\Gamma^{-1}y
    = \sum_{i=1}^{q}\mu_{i}^{j}P_{\mathcal{G}_i}\Gamma^{-1}y .
  \end{equation*}
  By \Cref{thm:qsggkb-properties}, we have $\hat{p}_{i}^{(k_b)}\in\mathcal{K}_{k_b}(\Gamma^{-1}M,\Gamma^{-1}y)=\mathrm{span}\{\widetilde{w}_i\}_{i=0}^{k_b-1}$, and
	\begin{equation*}
		(\widetilde{w}_{0},\dots,\widetilde{w}_{k_b-1}) = (P_{\mathcal{G}_1}\Gamma^{-1}y, \dots, P_{\mathcal{G}_q}\Gamma^{-1}y)
		\begin{pmatrix}
			1 & \mu_{1} & \cdots & \mu_{1}^{k_b-1} \\
			1 & \mu_{2} & \cdots & \mu_{2}^{k_b-1} \\
			\vdots & \vdots & \ddots & \vdots \\
			1 & \mu_{q} & \cdots & \mu_{q}^{k_b-1}
			\end{pmatrix} =: \widetilde{G}T_{k_b},
	\end{equation*} 
	where $\widetilde{G}=(P_{\mathcal{G}_1}\Gamma^{-1}y, \dots, P_{\mathcal{G}_q}\Gamma^{-1}y)$. Since the $q$-by-$q$ leading part of $T_{k_b}$ is a Vandermonde matrix with $\mu_i\neq\mu_j$, the rank of $T_{k_b}$ is at most $q$. Thus the rank of $\widetilde{W}:=(w_{0},\dots,w_{k_b-1})$ is at most $q$, which is achieved if $k_b=q$. This implies that the expansion of $\mathcal K_{k}(\Gamma^{-1}M,\Gamma^{-1}y)$ ends at $k=q$, hence $k_b =q$. Since $T_{k_b}=T_{q}$ is nonsingular, it holds that $\calR(\widetilde{W})=\calR(\widetilde{G})$, and we can write $\hat{p}_{i}^{(k_b)}$ as $\hat{p}_{i}^{(k_b)}=\widetilde{G}c\in\calG$ with a nonzero $c\in\mathbb{R}^{q}$. Now we immediately obtain $P_{\calG}\hat{p}_{i}^{(k_b)}=\hat{p}_{i}^{(k_b)}$, which is the desired result.
\end{proof}

\end{document}
